\section{Construct the loader without leaking a descriptor}
The first ten construction steps develop the loader. Nine named source regions hold
its implementation; the tenth step compares completed boundary cases. Fill one region
at a time. Its \texttt{BEGIN} and \texttt{END} comments identify the replacement span.
Keep the supplied prologue, checkpoints and restore paths. Checkpoint labels precede
ordinary \texttt{nop} instructions, so stopping before one shows the state produced
by the preceding region without executing that checkpoint instruction.

An unfinished entry returns development status four and count zero. Before open, an
unfinished region uses the cleanup path that owns no descriptor. After a successful
open, it uses the path that closes the retained descriptor first. After normal close,
it again uses the path that owns none. These are three points in one resource lifetime;
using the same close operation indiscriminately could close an invalid descriptor or
close the acquired one twice.

The supplied caller recognizes development status and prints
\texttt{Program is still under construction.} Its normal exit in this situation lets
you inspect an incomplete program; development status does not satisfy the completed
loader or decoder contract. Even after the loader is complete, that message remains
expected until the decoder entry is implemented. Inspect \texttt{loaded\_count} and
\texttt{last\_status} to distinguish progress in the two stages.

For each observation, prepare the named recipe from your current working source and
create its named fixture alongside the launched project. Direct-entry checks initialize
the same declared arguments and saved state; caller checks follow the actual supplied
main program. The draft recipe metadata supports these observations. It does not yet
implement an interactive trainer, remote progression control or enforced hint policy.
